\documentclass{article}

\usepackage{float}
\restylefloat{table}

\usepackage{booktabs}

\title{Team Productivity: Final\\\progname}

\author{\authname}

\date{}

\input{../Comments.tex}
\input{../Common.tex}

\begin{document}

\maketitle

This document summarizes the contributions of each team member for the final
demonstration and documentation.  The time period of interest is the time
between Rev 0 and the Final documentation; the contributions prior to Rev0 are
NOT included.

% \wss{Please do not delete sections or modify the table formats.  Standardization
% helps the instructors and TAs to review these documents. You may add sections to
% the report.}

% \wss{Remove the instructor guidelines/comments in the final version.}

% \wss{Fill in the requested information, like project url, TA name, supervisor name, etc.}

\section{Team Meeting Attendance}

% \wss{For each team member how many team meetings have they attended over the
% time period of interest.  This number should be determined from the meeting
% issues in the team's repo.  The first entry in the table should be the total
% number of team meetings held by the team.}

\begin{table}[H]
\centering
\begin{tabular}{ll}
\toprule
\textbf{Student} & \textbf{Meetings}\\
\midrule
Total & 2\\
Ethan & 1\\
Hussain & 2\\
Jeffrey & 1\\
Kevin & 1\\
Chengze & 2\\
\bottomrule
\end{tabular}
\end{table}

% \wss{If needed, an explanation for the counts can be provided here.}

One meeting was a user test meeting that only had Hussain and Chengze.

\section{Supervisor/Stakeholder Meeting Attendance}

% \wss{For each team member how many supervisor/stakeholder team meetings have
% they attended over the time period of interest.  This number should be determined
% from the supervisor meeting issues in the team's repo.  The first entry in the
% table should be the total number of supervisor and team meetings held by the
% team.  If there is no supervisor, there will usually be meetings with
% stakeholders (potential users) that can serve a similar purpose.}

\noindent \textbf{Supervisor's Name: } [N/A]

\begin{table}[H]
\centering
\begin{tabular}{ll}
\toprule
\textbf{Student} & \textbf{Meetings}\\
\midrule
Total & 0\\
Ethan & 0\\
Hussain & 0\\
Jeffery & 0\\
Chengze & 0\\
Kevin & 0\\
\bottomrule
\end{tabular}
\end{table}

\wss{If needed, an explanation for the counts can be provided here.}

\section{Lecture Attendance}

% \wss{For each team member how many lectures have they attended over the time
% period of interest.  This number should be determined from the lecture issues in
% the team's repo. You can find the number of lectures in the time period of
% interest by looking at the
% \href{https://calendar.google.com/calendar/u/0/embed?src=rnboqiaki1k2la7rpu3bn0um58@group.calendar.google.com&ctz=America/Toronto}
% {Google calendar} for the capstone course. The first entry in the
% table should be the total number of lectures held for the class during the period of interest.}

% \wss{NOTE: There will be approximately 1 lecture between the Rev0 and Rev1
% demos}

\begin{table}[H]
\centering
\begin{tabular}{ll}
\toprule
\textbf{Student} & \textbf{Lectures}\\
\midrule
Total & 1\\
Ethan & 1\\
Hussain & 1\\
Jeffery & 1\\
Chengze & 1\\
Kevin & 0\\
\bottomrule
\end{tabular}
\end{table}

\wss{If needed, an explanation for the lecture attendance can be provided here.}

\section{TA Document Discussion Attendance}

\noindent \textbf{TA's Name: } [fill in this information]

\begin{table}[H]
\centering
\begin{tabular}{ll}
\toprule
\textbf{Student} & \textbf{Meetings}\\
\midrule
Total & 1\\
Ethan & 1\\
Hussain & 1\\
Jeffrey & 1\\
Kevin & 1\\
Chengze & 1\\
\bottomrule
\end{tabular}
\end{table}


\section{Commits, Lines Added and Lines Deleted, Pull Requests}

\begin{table}[H]
\centering
\begin{tabular}{lllllll}
\toprule
\textbf{Student} & \textbf{Commits} & \textbf{Percent} & \textbf{Lines Added} &
\textbf{Lines Deleted} & \textbf{Pull Requests}\\
\midrule
Total & 123 & 100\% \\
Ethan & 44 & 35.7\% & +2652 & -623 & 18\\
Hussain & 36 & 29.3\% & +1336 & -230 & 12\\
Jeffrey & 15 & 12.2\% & +909 & -187 & 8\\
Kevin & 13 & 10.5\% & +209 & -24 & 4\\
Chengze & 15 & 12.2\% & +765 & -49 & 10\\
\bottomrule
\end{tabular}
\end{table}

\section{Issue Tracker}

% \wss{For each team member how many issues have they authored (including open and
% closed issues (O+C)) and how many have they been assigned (only counting closed
% issues (C only)) over the time period of interest.}

\begin{table}[H]
\centering
\begin{tabular}{lll}
\toprule
\textbf{Student} & \textbf{Authored (O+C)} & \textbf{Assigned (C only)}\\
\midrule
Ethan & 7 & 6 \\
Hussain & 11 & 5 \\
Jeffrey & 2 & 3 \\
Kevin & 1 & 2 \\
Chengze & 0 & 4 \\
\bottomrule
\end{tabular}
\end{table}

% \wss{If needed, an explanation for the counts can be provided here.}

\section{Team Charter Trigger Items}

\wss{Provide a summary of the quantified triggers identified in the team's
charter.}

\wss{Provide a list of any violations of the triggers.  If the team wishes, the
violations can be summarized on aggregate, instead of naming specific team
members.}

\wss{Provide a plan to address the violations.  This could include revising the
triggers, if they are found to be too weak, strong or ambiguous.}

\section{Additional Productivity Metrics}

Our team does not have any additional metrics of
productivity aside from the ones listed in this report.

\end{document}
